\documentclass[12pt]{article}
\usepackage[utf8]{inputenc}
\usepackage{amsmath,amssymb}
\usepackage{geometry}
\usepackage{tikz}
\usepackage{caption}
\geometry{a4paper,margin=1in}
\setcounter{secnumdepth}{0}

\title{From a Single Jump to a Point Mass: A Story of Heaviside and Dirac}
\author{Haiyi Li}
\date{}

\begin{document}
\maketitle

\section{Opening: Why Revisit a ``Simple'' Function?}

It seems we may have overlooked some special properties of very simple functions
that we use all the time.
In my study of integral transforms, I was struck by a basic observation:
a step function with only one jump point already carries unexpectedly rich content
in distribution theory and measure theory \cite{AndrewsShivamoggi1999}.
This note tells that story, from a jump discontinuity to a point mass.

\section{Act I: The Step Function and Its First Shape}

Define
\[
 H(t-a)=
 \begin{cases}
 0, & t<a,\\
 1, & t>a.
 \end{cases}
\]
This function has a unit jump at $t=a$ and acts like a switch.
The value at $t=a$ can be chosen by convention (for instance $H(0)=\tfrac12$),
but that choice does not affect the distributional derivative.

\begin{figure}[h]
  \centering
  \begin{tikzpicture}[>=stealth,scale=1]
    % axes
    \draw[->] (-2,0) -- (4,0) node[right] {$t$};
    \draw[->] (0,-0.5) -- (0,2.0) node[above] {$H(t-a)$};
    % step
    \draw (-2,0) -- (1,0);
    \draw[dashed] (1,0) -- (1,1.3);
    \draw (1,1.3) -- (4,1.3);
    % labels
    \node[below] at (1,0) {$a$};
    \node[left] at (0,1.3) {$1$};
  \end{tikzpicture}
  \caption{The Heaviside unit step}
\end{figure}

A rectangle over $[a,b]$ can be written as
\[
 f(t)=H(t-a)-H(t-b).
\]
So even at this stage, a jump already builds geometric structure.

\begin{figure}[h]
  \centering
  \begin{tikzpicture}[>=stealth,scale=1]
    % axes
    \draw[->] (-2,0) -- (4,0) node[right] {$t$};
    \draw[->] (0,-0.5) -- (0,2.0);
    % rectangle
    \draw (-2,0) -- (-0.5,0);
    \draw[dashed] (-0.5,0) -- (-0.5,1.3);
    \draw (-0.5,1.3) -- (2.2,1.3);
    \draw[dashed] (2.2,1.3) -- (2.2,0);
    \draw (2.2,0) -- (4,0);
    % labels
    \node[below] at (-0.5,0) {$a$};
    \node[below] at (2.2,0) {$b$};
  \end{tikzpicture}
  \caption{Rectangle function $f(t)=H(t-a)-H(t-b)$}
\end{figure}

\section{Act II: Squeezing Area into a Point}

Consider a rectangular pulse of width $2\varepsilon$ and height $\frac{1}{2\varepsilon}$:
\[
 d_a(t)=\frac{1}{2\varepsilon}\Bigl[H(t-a+\varepsilon)-H(t-a-\varepsilon)\Bigr],
\]
with unit area. Letting $\varepsilon\to 0$ yields the idealized impulse (Dirac $\delta$):
\[
 \delta(t-a)=\lim_{\varepsilon\to 0} d_a(t),
 \qquad
 \int_{-\infty}^{\infty}\delta(t-a)\,dt=1.
\]

\begin{figure}[h]
  \centering
  \begin{tikzpicture}[>=stealth,scale=1]
    % axes
    \draw[->] (-2,0) -- (4,0) node[right] {$t$};
    \draw[->] (0,-0.5) -- (0,2.0) node[above] {$d_a(t)$};
    % pulse
    \draw (-2,0) -- (0.2,0);
    \draw[dashed] (0.2,0) -- (0.2,1.3);
    \draw (0.2,1.3) -- (1.8,1.3);
    \draw[dashed] (1.8,1.3) -- (1.8,0);
    \draw (1.8,0) -- (4,0);
    % labels
    \node[below] at (1.0,0) {$a$};
    \node[below] at (0.2,0) {$a-\varepsilon$};
    \node[below] at (1.8,0) {$a+\varepsilon$};
    \node[left] at (0,1.3) {$\frac{1}{2\varepsilon}$};
  \end{tikzpicture}
  \caption{Rectangular approximation of the impulse}
\end{figure}

\section{Act III: The Distribution Appears}

For a suitable test function $f$,
\[
 \int_{-\infty}^{\infty} \delta(t-a) f(t)\,dt = f(a).
\]
This sifting identity shows why $\delta$ is not an ordinary function:
it is meaningful through its action inside integrals, i.e., as a distribution.

Now the key step in the story:
for a test function $\varphi\in C_c^\infty(\mathbb{R})$,
\[
 \langle (H(t-a))',\varphi\rangle
 =-\langle H(t-a),\varphi'\rangle
 =-\int_a^\infty \varphi'(t)\,dt
 =\varphi(a)
 =\langle \delta(t-a),\varphi\rangle.
\]
Hence, in the distributional sense,
\[
 (H(t-a))'=\delta(t-a).
\]
So the derivative of the jump is a point impulse.

\section{Act IV: Measure-Theoretic Meaning (Core Viewpoint)}

In measure language,
for a point $a$, the Dirac measure $\delta_a$ assigns mass $1$ to any set containing
$a$ and $0$ otherwise, so the sifting property becomes
\[
 \int f(t)\,d\delta_a(t)=f(a).
\]
The Heaviside function is discontinuous, but as a distribution (or equivalently, in the
Stieltjes integral sense) its derivative is a measure concentrated at the jump point.
Thus the statement $(H(t-a))'=\delta(t-a)$ expresses that the distributional derivative
of a jump is a point mass.

\paragraph{Lebesgue measure (context).}
The Lebesgue measure on $\mathbb{R}$ is the standard measure used to quantify the
``size'' of sets. For an interval $[a,b]$ it equals the length,
\[
 \mu([a,b])=b-a.
\]
It is a regular measure on the Borel sets and is countably additive: for any countable
collection of measurable sets $\{A_i\}$,
\[
 \mu\Bigl(\bigcup_i A_i\Bigr)=\sum_i \mu(A_i).
\]
Lebesgue measure applies to many kinds of sets (for instance, dense sets and the set of
irrationals), but it assigns zero measure to every singleton:
\[
 \mu(\{a\})=0 \quad \text{for all } a\in\mathbb{R}.
\]

\paragraph{Dirac vs.\ Lebesgue (measure viewpoint).}
The Dirac $\delta$ is not an ordinary function; it is a distribution that corresponds
to the \emph{Dirac measure} $\delta_a$, a point mass at $a$. In contrast, Lebesgue
measure describes the size of extended sets. A useful comparison:
\begin{itemize}
  \item \textbf{Support.} Lebesgue measure has support $\mathbb{R}$, while $\delta_a$
  is supported only at the single point $t=a$.
  \item \textbf{Values.} Lebesgue measure gives $\mu(\{a\})=0$ and
  $\mu([a,b])=b-a$. Dirac measure gives $\delta_a(A)=1$ if $a\in A$ and $0$ otherwise.
  \item \textbf{Additivity and regularity.} Both are genuine measures: $\mu$ and
  $\delta_a$ are countably additive and regular on Borel sets. The Dirac measure is
  ``extreme'' only in the sense that all of its mass is concentrated at a single point.
\end{itemize}

\section{Act V: Consequences and Proofs (Preserving the Toolkit)}

\paragraph{Consequences and proofs.}
We use the distributional identity $f(t)\delta(t-a)=f(a)\delta(t-a)$.

\textit{(a) $t\,\delta(t)=0$.}
Let $f(t)=t$. Then $t\delta(t)=f(t)\delta(t)=f(0)\delta(t)=0\cdot\delta(t)=0$.

\textit{(b) $\delta(a t)=\frac{1}{|a|}\delta(t)$ for $a\neq 0$.}
For any test function $\varphi$,
\[
 \int_{-\infty}^{\infty}\delta(a t)\varphi(t)\,dt
 =\int_{-\infty}^{\infty}\delta(u)\varphi(u/a)\,\frac{du}{|a|}
 =\frac{1}{|a|}\varphi(0)
 =\int_{-\infty}^{\infty}\frac{1}{|a|}\delta(t)\varphi(t)\,dt.
\]
Thus $\delta(a t)=\frac{1}{|a|}\delta(t)$.

\textit{(c) $\delta(t-a)=\delta(a-t)$.}
Use (b) with $a=-1$:
\[
\delta(a-t)=\delta(-(t-a))=\delta(-1\cdot(t-a))=\delta(t-a).
\]

\textit{(d) If $g$ is monotone, $C^1$ near $a$, $g(a)=0$, and $g'(a)\neq 0$, then
$\delta(g(t))=\frac{1}{|g'(a)|}\delta(t-a)$.}
For a test function $\varphi$, substitute $u=g(t)$ (monotonicity gives a valid change
of variables). Then
\[
 \int \delta(g(t))\varphi(t)\,dt
 =\int \delta(u)\,\varphi(g^{-1}(u))\,\frac{du}{|g'(g^{-1}(u))|}
 =\frac{\varphi(a)}{|g'(a)|}.
\]
Hence $\delta(g(t))=\frac{1}{|g'(a)|}\delta(t-a)$.

\textit{(e) From (d), $\delta(a t-b)=\frac{1}{|a|}\delta(t-b/a)$.}
Let $g(t)=a t-b$ so $g(a_0)=0$ at $a_0=b/a$ and $g'(t)=a$.

\section{Closing Remark}

The central message is simple but deep:
in classical calculus, jumps look singular; in measure and distribution language,
they become precise objects. The step function does not merely ``jump'' at a point;
its derivative records that jump as a point mass.

\bibliographystyle{plain}
\bibliography{references}

\end{document}
